% Experiment report. Fixed sections: see experiments/TEMPLATE/report.tex and experiments/README.md.
% Build: npm run experiment:report -- 09-reserve
\documentclass{../polyexp}
\begin{document}
% generated by scripts/experiment-report.ts from experiment.json and models/<run>/ - do not edit
\providecommand{\ExpHeader}{}\renewcommand{\ExpHeader}{%
{\sffamily\footnotesize\color{expgrey}EXPERIMENT 9 \textperiodcentered\ MECHANISM \hfill 2026-11-24}\par
{\LARGE\bfseries\sffamily Saving stars for a better purchase}\hfill\verdictbox{expred}{NOT KEPT}\par\smallskip
{\itshape Does a one-step greedy bot gain from skipping a cheap purchase that would leave it unable to afford a better one next turn?}\par\medskip
\begin{tabular}{@{}p{22mm}p{\dimexpr\linewidth-26mm}@{}}
\fact{Change}{A purchase that leaves next turn's stars short of the best purchase then affordable, and scores below it, loses reserve times that purchase's score.}
\fact{Knobs}{\knob{reserve} $\in$ \{0.15, 0.3, 0.6, 1\}, best 0.15}
\fact{Base}{\model{greedy-full-front}}
\fact{Method}{Strength sweep, 200 games per point against each panel model, final test 400 games; metric: points share}
\fact{Setups}{4 players, 16$\times$16, drylands, turn cap 150}
\fact{Record}{issue \issue{41}, pull request \issue{126}; status: done}
\end{tabular}\par\medskip
\begin{center}\begin{tikzpicture}[node distance=5mm]
    \node[flowbox] (n0) {{\sffamily\bfseries\scriptsize\color{expgrey}BASE}\\[1pt]\model{full-front}};
    \node[flowbox,right=5mm of n0] (n1) {{\sffamily\bfseries\scriptsize\color{expgrey}CHANGE}\\[1pt]\knob{reserve}\\{\scriptsize 0.15, 0.3, 0.6, 1}};
    \node[flowbox,right=5mm of n1] (n2) {{\sffamily\bfseries\scriptsize\color{expgrey}MEASURED}\\[1pt]Sweep, 200 games per point\\{\scriptsize 1 setup, metric: points share}};
    \node[flowbox,right=5mm of n2] (n3) {{\sffamily\bfseries\scriptsize\color{expgrey}RESULT}\\[1pt]Points vs base, final test\\{\scriptsize 50.0 $\rightarrow$ 50.0\,\%}\\\textbf{+0.0\pp}};
    \node[flowbox,right=5mm of n3] (n4) {{\sffamily\bfseries\scriptsize\color{expgrey}VERDICT}\\[1pt]\textbf{\color{expred}NOT KEPT}};
    \draw[flowarrow] (n0) -- (n1); \draw[flowarrow] (n1) -- (n2); \draw[flowarrow] (n2) -- (n3); \draw[flowarrow] (n3) -- (n4);
\end{tikzpicture}\end{center}
\begin{center}\sffamily\small\begin{tabular}{@{}lrrr@{}}\toprule
Measure & Base & New & Change\\\midrule
Points vs base, final test & 50.0\,\% & 50.0\,\% & +0.0\\
Points share, 6-model panel, sweep & 70.3\,\% & 70.3\,\% & +0.0\\
Decisions changed per 1,000 & 0.0 & 2.0 & +2.0\\
\bottomrule\end{tabular}\end{center}
}
\providecommand{\ExpResults}{}\renewcommand{\ExpResults}{%
\subsection{Sweep, 4 players {\normalfont\small(\model{greedy-full-reserve})}}
\begin{center}\sffamily\small\begin{tabular}{@{}lrrrrrr@{}}\toprule
Point & Finish (\%) & Points (\%) & Won (\%) & Conquest (\%) & Turns & vs base: points (\%)\\\midrule
base & -- & 67.36 & 80.7 & -- & -- & 50.0\\
reserve=0.15 & -- & 67.40 & 80.8 & -- & -- & 50.0\\
reserve=0.3 & -- & 67.40 & 80.8 & -- & -- & 50.0\\
reserve=0.6 & -- & 67.36 & 80.9 & -- & -- & 48.8\\
reserve=1 & -- & 67.31 & 81.2 & -- & -- & 49.1\\
\bottomrule\end{tabular}\\[3pt]{\footnotesize\color{expgrey}Panel means; 200 games against each of 7 panel models per point.}\end{center}
\begin{center}\begin{tikzpicture}\begin{axis}[expplot,height=4.6cm,xtick={0,1,2,3,4},xticklabels={{base},{0.15},{0.3},{0.6},{1}},
  xlabel={\texttt{reserve}},ylabel={Points share, panel mean (\%)},enlarge x limits=0.12,enlarge y limits=0.25]
\addplot[dashed,expgrey,line width=0.7pt,forget plot] coordinates {(0,67.357) (4,67.357)};
\addplot[only marks,mark=*,mark size=2.6pt,expblue] coordinates { (1,67.400) (2,67.400) (3,67.357) (4,67.314)};
\addplot[only marks,mark=*,mark size=2.6pt,expgrey] coordinates {(0,67.357)};
\end{axis}\end{tikzpicture}\\[-2pt]{\footnotesize\color{expgrey}Dashed: the base. Grey: base; blue: sweep points.}\end{center}
\begin{center}\sffamily\small\begin{tabular}{@{}lrr@{}}\toprule
Opponent & Won (\%) & Points (\%)\\\midrule
\model{v1} & 94 $\pm$ 2 & 75.6\\
\model{hill-v1} & 91 $\pm$ 3 & 72.3\\
\model{hill-v2} & 95 $\pm$ 2 & 73.7\\
\model{hill-v3} & 91 $\pm$ 3 & 70.3\\
\model{hill-v4} & 87 $\pm$ 3 & 68.6\\
\model{full-v1} & 58 $\pm$ 5 & 54.8\\
\model{full-front} & 50 $\pm$ 5 & 50.0\\
\midrule
\textbf{Mean without \model{full-front}} & 85.9 & 69.2\\
\bottomrule\end{tabular}\\[3pt]{\footnotesize\color{expgrey}Final test of reserve=0.15: 400 unseen games against each.}\end{center}
}

\ExpHeader

\section{Question and hypothesis}
The greedy bot spends stars on the best affordable action every step. The Autonomous Systems course calls this kind of one-step choice myopic and motivates planning as looking ahead and delaying action (AS 2.0 Planning as Search, PDF page 15, slide 8/51; technique AS-05, issue \#41). \textbf{Hypothesis:} skipping a cheap purchase that would block a clearly better one next turn raises the points share against the panel by about one point and wins the head-to-head against the base.

\section{Mechanism}
Each turn the bot finds the best purchase it cannot afford now but can afford next turn (stars set to next turn's income, cost in $(\mathit{stars}, \mathit{next}]$), with score $R$ and cost $C$. A purchase with cost $c>0$ and score $0<s<R$ for which $\mathit{next}-c<C$ scores $s-\knob{reserve}\cdot R$ instead. Default 0 (off). Code: \texttt{purchase\_cost}, \texttt{reserve\_reference} and the saving branch in \texttt{full\_best\_of}, \texttt{engine/core/src/bot\_full.rs}.

\section{Setup}
Values were sized so that at 1.0 a blocked purchase loses its full reference score, enough to push any cheaper purchase below 0. This experiment ran before the finish metric (\#124), so it uses the points share. Base: \model{greedy-full-front}, the champion at the time.

\section{Results}
\ExpResults

All four points lie within 0.2 points of the base on the 6-model panel (70.25\,\%), inside the noise of about $\pm$1.2. The final test of the best point (0.15) against the base is an exact tie: 50.0\,\% points, 50\,\% won ($\pm$5).

\section{Behaviour}
At \knob{reserve} 1.0 the mechanism changed 4 of 1,879 bot decisions (about 2 per 1,000) in the first 12 turns of generated 4-player games. The knob saturates at 1.0: larger values change nothing more.

\section{Interpretation}
The prediction failed because the condition almost never holds: a better purchase affordable next turn but not now, blocked by a cheaper one now. The idea is sound, but a one-turn lookahead on a single reference purchase is too narrow to change play. A plan over several turns is the DP of experiment 10.

\todo{jakob}{Your reading: agree, disagree, what is missing?}

\section{Threats to validity}
\begin{itemize}
\item The behaviour count comes from the first 12 turns only; later turns, with more stars, may trigger it more often.
\item The reference is one purchase. A plan that weighs several future purchases might act often enough to show an effect.
\end{itemize}

\section{Verdict and next step}
\textbf{Not kept}: the head-to-head is a tie and the panel moves less than the noise. Merged off by default (PR \#126) as the record. Champion stays \model{greedy-full-front}. Next: the multi-turn version, an economy plan by dynamic programming (experiment 10, \#68).

\section{Notes for the paper}
Draft: a negative result worth one line. A myopia patch that looks one turn ahead at one purchase almost never fires (2 per 1,000 decisions); the gain needs a real multi-step plan (experiment 10).

\todo{jakob}{Is this worth a place in the paper? Where?}

\end{document}
